\documentclass[10pt]{article}
% Shared typographic and algebraic source.
\usepackage[letterpaper,margin=0.86in]{geometry}
\usepackage[T1]{fontenc}
\usepackage[utf8]{inputenc}
\usepackage{lmodern}
\usepackage{microtype}
\usepackage{amsmath,amssymb,amsthm,mathtools,mathrsfs}
\usepackage{booktabs,array,tabularx,longtable}
\usepackage{enumitem}
\usepackage[dvipsnames]{xcolor}
\usepackage{fancyhdr,titlesec}
\usepackage[hidelinks]{hyperref}

\definecolor{ink}{HTML}{1F2A33}
\definecolor{teal}{HTML}{226A73}
\definecolor{muted}{HTML}{65747E}
\definecolor{line}{HTML}{C5D1D3}
\definecolor{wash}{HTML}{F3F7F6}

\setlength{\parindent}{0pt}
\setlength{\parskip}{0.46em}
\setlength{\emergencystretch}{2em}
\allowdisplaybreaks[2]
\renewcommand{\arraystretch}{1.16}
\setlist{nosep,leftmargin=1.5em}

\pagestyle{fancy}
\fancyhf{}
\fancyhead[L]{\footnotesize\sffamily\color{muted}COPPALINI ARCHITECTURE}
\fancyhead[R]{\footnotesize\sffamily\color{muted}PHYSICAL REALITY INSIDE CONSTRAINT}
\fancyfoot[C]{\footnotesize\sffamily\color{muted}\thepage}
\renewcommand{\headrulewidth}{0.35pt}
\renewcommand{\headrule}{\hbox to\headwidth{\color{line}\leaders\hrule height \headrulewidth\hfill}}
\setlength{\headheight}{19pt}

\titleformat{\section}{\Large\bfseries\color{ink}}{\thesection.}{0.55em}{}
\titleformat{\subsection}{\large\bfseries\color{teal}}{\thesubsection.}{0.5em}{}
\titlespacing*{\section}{0pt}{1.42em}{0.48em}
\titlespacing*{\subsection}{0pt}{1.02em}{0.32em}

\newtheorem{theorem}{Theorem}
\newtheorem{proposition}[theorem]{Proposition}
\newtheorem{lemma}[theorem]{Lemma}
\theoremstyle{definition}
\newtheorem{definition}[theorem]{Definition}
\theoremstyle{remark}
\newtheorem{remark}[theorem]{Boundary}

\newcommand{\RR}{\mathbb R}
\newcommand{\CC}{\mathbb C}
\newcommand{\QQ}{\mathbb Q}
\newcommand{\ZZ}{\mathbb Z}
\newcommand{\Id}{I}
\newcommand{\tr}{\operatorname{tr}}
\newcommand{\im}{\operatorname{im}}
\newcommand{\kerop}{\operatorname{ker}}
\newcommand{\Span}{\operatorname{span}}
\newcommand{\diag}{\operatorname{diag}}
\newcommand{\snoc}{\operatorname{snoc}}
\newcommand{\Admit}{\operatorname{Admit}}
\newcommand{\Gate}{\operatorname{Gate}}
\newcommand{\Typed}{\operatorname{Typed}}
\newcommand{\Future}{\operatorname{future}}
\newcommand{\Scope}{\operatorname{sc}}
\newcommand{\SilScope}{\mathsf{Silence}_{\rm scope}}
\newcommand{\SilFit}{\mathsf{Silence}_{\rm fit}}

\newcommand{\papertitle}[2]{%
  \begin{center}
  {\color{teal}\rule{\textwidth}{1.35pt}}\\[0.38in]
  {\footnotesize\sffamily\bfseries\color{teal} COPPALINI ARCHITECTURE}\\[0.16in]
  {\fontsize{31}{36}\selectfont\bfseries\color{ink}#1\par}
  \vspace{0.15in}
  {\Large\color{teal}#2\par}
  \vspace{0.22in}
  {\color{line}\rule{\textwidth}{0.55pt}}\\[0.12in]
  {\normalsize Anthony Vito Coppa}\par
  {\small 9 November 2026}\par
  \vspace{0.30in}
  \end{center}
}

\newcommand{\sharedroot}{%
Let
\[
Q=\begin{pmatrix}1&1\\1&0\end{pmatrix},\qquad
\Omega=\begin{pmatrix}0&1\\-1&0\end{pmatrix},\qquad
C=-\Omega,
\]
\[
G=[Q,\Omega]=\begin{pmatrix}-2&1\\1&2\end{pmatrix},\qquad
H=\frac{G}{\sqrt5},\qquad
S=\frac{2Q-I}{\sqrt5}.
\]
The common exact identities are
\[
Q^2=Q+I,\qquad C^2=-I,\qquad G^2=5I,\qquad H^2=S^2=I,
\]
\[
H=CS,\qquad SG=\sqrt5\Omega=-GS,\qquad
L_s:=e^{sS},\qquad
L_s^{\mathsf T}GL_s=G,\qquad HL_sH=L_{-s}.
\]
}

\newcommand{\sharedoccurrence}[2]{%
A physical field has an interface
\[
\begin{aligned}
\mathfrak I_J:=(&J,\Gamma_J,\operatorname{Src}_J;
A_J,X_J,Y_J,T_J,\Sigma_J,\operatorname{Con}_J,\\
&\rho_J,O_J,K_J,B_J,\theta_J),
\qquad
\rho_J:X_J\to Y_J,\quad
O_J:X_J\to\mathcal O_J,\quad
\theta_J:X_J\to\{0,1\}.
\end{aligned}
\tag{#1}
\]
Its occurrence type is
\[
\begin{aligned}
\mathsf{Occ}_J:=\{&
(\mathfrak I_J;\Theta,x,\rho_J(x),\kappa,
\boldsymbol\varepsilon,\theta_J(x)):\\
&\Theta\in A_J,\ x\in X_J,\ \kappa\in K_J,\
\boldsymbol\varepsilon\in\RR_{\ge0}^{3}\}.
\end{aligned}
\tag{#2}
\]
An indexed occurrence $\mathsf{Occ}_{J,n}$ is an element of
$\mathsf{Occ}_J$.
}

\newcommand{\sharedmatching}{%
Let
\[
\mathcal A_8:=\{\infty\}\sqcup\ZZ_7,
\qquad
(\infty,0,1,2,3,4,5,6)
=(\mathsf{Gr},\mathsf{Li},\mathsf{Me},\mathsf{Ti},
\mathsf{Hz},\mathsf{Si},\mathsf{Sc},\mathsf{Mo}).
\]
For $r\in\ZZ_7$, set
\[
\mathsf M_r=\{\{\infty,r\}\}\cup
\{\{r+j,r-j\}:j=1,2,3\}.
\]
Then $E(K_8)=\bigsqcup_{r\in\ZZ_7}\mathsf M_r$, and the canonical matching is
\[
\mathsf M_0=\{\mathsf{Gravity\! -\! Light},\;
\mathsf{Memory\! -\! Motion},\;
\mathsf{Time\! -\! Scope},\;
\mathsf{Horizon\! -\! Silence}\}.
\]
}

\hypersetup{pdftitle={Time Scope: Local Phase and Projective Return},pdfauthor={Anthony Vito Coppa},pdfsubject={Physical phase, projective scope, and half-turn return}}
\begin{document}
\papertitle{Time Scope}{Local phase and projective return}

\begin{abstract}
Time carries local phase through physical distance and duration. Scope reads
the direction of a state through its projective identity. A half-turn changes
every nonzero local state while leaving its projective scope unchanged. The
mirror reverses phase. Linear polarization and a two-level quantum state carry
the same double-angle relation: local phase advances by $\alpha$ while the
scope read advances by $2\alpha$. The field therefore separates local change,
projective return, measurement direction, and time-addressed motion.
\end{abstract}

\section{Time and scope}

For a wave of angular frequency $\omega$ and wave covector $\mathbf k$, phase
difference across a spacetime interval is
\[
\Delta\alpha=\omega\,\Delta t-\mathbf k\cdot\Delta\mathbf x.
\tag{TS.0}
\]
For light in vacuum,
\[
\omega=c|\mathbf k|,
\qquad
\lambda_{\rm wave}=\frac{2\pi}{|\mathbf k|},
\qquad
c=\lambda_{\rm wave}f.
\tag{TS.0a}
\]
For a material trajectory with velocity $\mathbf v$,
\[
\Delta\mathbf x=\mathbf v\,\Delta t,
\qquad
\Delta\alpha=(\omega-\mathbf k\cdot\mathbf v)\Delta t.
\tag{TS.0b}
\]
Thus duration, distance, wavelength, frequency, and light speed determine one
physical phase relation. The local phase carrier is
$V^\times=\RR^2\setminus\{0\}$. Scope reads the unoriented line through its
state.

The half-turn relation is
\[
\Scope(U_\pi z)=\Scope(z),\qquad U_\pi z\ne z
\quad(z\in V^\times).
\tag{TS.1}
\]
It is a return of scope through local difference.

\section{Phase carrier and mirror}

\sharedroot

Define the phase transport
\[
U_\alpha=e^{\alpha C}
=\begin{pmatrix}\cos\alpha&-\sin\alpha\\
\sin\alpha&\cos\alpha\end{pmatrix},
\qquad \alpha\in\RR/2\pi\ZZ.
\tag{TS.2}
\]

\begin{proposition}[Phase and mirror]
For all $\alpha$,
\[
U_{\alpha+\beta}=U_\alpha U_\beta,\qquad
HU_\alpha H=U_{-\alpha},
\qquad U_\pi=-I.
\]
\end{proposition}

\begin{proof}
The composition law is the one-parameter exponential law. Since $HC=-CH$,
$H e^{\alpha C}H=e^{-\alpha C}$. Finally
$e^{\pi C}=I\cos\pi+C\sin\pi=-I$.
\end{proof}

The split form supplies two null lines
\[
\ell_\pm=\RR\left(1,\frac{-1\pm\sqrt5}{2}\right)^{\mathsf T},
\qquad \ell_\pm^{\mathsf T}G\ell_\pm=0.
\tag{TS.3}
\]
The mirror and quarter-turn exchange these two projective directions.

\section{Scope and the half-turn receipt}

\begin{definition}[Scope reader]
The scope reader is the projectivization
\[
\Scope:V^\times\longrightarrow\RR P^1,\qquad \Scope(z)=[z].
\tag{TS.4}
\]
It forgets nonzero scalar magnitude and sign while preserving the line through
the origin.
\end{definition}

Define the scope-closure selector
\[
\Phi:\RR/2\pi\ZZ\longrightarrow\mathcal P(V^\times),
\qquad
\Phi(\alpha):=
\{z\in V^\times:\Scope(U_\alpha z)=\Scope(z)\},
\qquad
\SilScope:=V^\times.
\tag{TS.4a}
\]

\begin{theorem}[Half-turn scope return]
For every $z\in V^\times$,
\[
U_\pi z=-z\ne z,\qquad \Scope(U_\pi z)=\Scope(z).
\tag{TS.5}
\]
Thus the half-turn receipt is
\[
\Phi(\pi)=\SilScope.
\tag{TS.6}
\]
\end{theorem}

\begin{proof}
The first equality is the preceding proposition. A nonzero vector and its
negative span the same one-dimensional real subspace, so their projective
readings are equal. A nonzero real vector cannot equal its negative.
\end{proof}

The projective action of $H$ is well typed:
\[
\overline H:\RR P^1\to\RR P^1,\qquad \overline H([z])=[Hz].
\tag{TS.7}
\]
It is a bijection and supplies scope covariance,
\[
\Scope(Hz)=\overline H(\Scope(z)).
\tag{TS.8}
\]

For an explicit circle coordinate, define
\[
\chi([x:y])=
\left(\frac{x^2-y^2}{x^2+y^2},\frac{2xy}{x^2+y^2}\right)\in S^1.
\tag{TS.9}
\]
If $z=(\cos\alpha,\sin\alpha)$, then
\[
\chi(\Scope(z))=(\cos2\alpha,\sin2\alpha).
\tag{TS.10}
\]
The doubled angle is the scope coordinate of the two-to-one phase relation.


\section{Ordered transport and memory return}

Let $(A,\preceq)$ be an address order with initial address $a_0$.
Let $\phi:A_{\succeq a_0}\to\RR/2\pi\ZZ$ be its phase coordinate. Define
\[
W_{b\leftarrow a}:=U_{\phi(b)-\phi(a)},
\qquad a\preceq b.
\tag{TS.11}
\]

\begin{proposition}[Functorial transport]
For $a\preceq b\preceq c$,
\[
W_{c\leftarrow b}W_{b\leftarrow a}=W_{c\leftarrow a},
\qquad W_{a\leftarrow a}=I.
\tag{TS.12}
\]
\end{proposition}

For $\psi_0\in V^\times$, the trajectory is
\[
\Psi_{a_0,\psi_0}(a):=W_{a\leftarrow a_0}\psi_0,
\tag{TS.13}
\]
and its trajectory-relative displacement is
\[
\delta_{\rm traj}(a):=\Psi_{a_0,\psi_0}(a)-\psi_0.
\tag{TS.14}
\]
The displacement is measured from the initial state at $a_0$.

\sharedoccurrence{TS.14a}{TS.14b}

For $J=\mathsf{TS}$, take $\Theta_n=a_n$,
$x_n=\Psi_{a_0,\psi_0}(a_n)$, and $y_n=\Scope(x_n)$. The time--scope
projection is
\[
\operatorname{pr}_{\rm TS}(\mathsf{Occ}_{\mathsf{TS},n})
=(\Theta_n,x_n,y_n,\delta_{\rm traj}(a_n),
\boldsymbol\varepsilon_n).
\tag{TS.14c}
\]

\begin{definition}[Memory return]
For a reader $\rho:V^\times\to Y$, define
\[
\operatorname{MemRet}_{\rho}(a,b;\psi_0)
\Longleftrightarrow
\begin{cases}
a\prec b,\\
\rho(\Psi(a))=\rho(\Psi(b)),\\
\exists m\ (a\prec m\prec b\ \land\ \rho(\Psi(m))\ne\rho(\Psi(a))).
\end{cases}
\tag{TS.15}
\]
\end{definition}

The intermediate address witnesses departure between the equal endpoint reads.
For the scope reader, a half-turn with an intermediate quarter-turn gives
memory return:
\[
\phi(b)-\phi(a)=\pi,\qquad
\phi(m)-\phi(a)=\frac\pi2
\quad\Longrightarrow\quad
\operatorname{MemRet}_{\Scope}(a,b;\psi_0).
\tag{TS.16}
\]

\section{Mirror transfer}

Define mirrored phase data by $\phi^\mu=-\phi$ and mirrored trajectory by
\[
\Psi^\mu(a):=H\Psi(a).
\tag{TS.17}
\]
Then $\Psi^\mu$ has transport
$W^\mu_{b\leftarrow a}=U_{-(\phi(b)-\phi(a))}$.

\begin{theorem}[Mirror transfer of memory return]
Suppose a reader $\rho$ has an injective mirror covariance
\[
\rho(Hz)=\overline H_\rho(\rho(z))
\tag{TS.18}
\]
for an injective map $\overline H_\rho$ on its image. Then
\[
\operatorname{MemRet}_{\rho}(a,b;\psi_0)
\Longleftrightarrow
\operatorname{MemRet}_{\rho}(a,b;H\psi_0)\ \text{on mirrored data}.
\tag{TS.19}
\]
\end{theorem}

\begin{proof}
The mirror relation gives $\Psi^\mu(x)=H\Psi(x)$ at every address $x$.
Covariance transfers endpoint equality. Injectivity transfers the strict
inequality at the intermediate departure address. The converse follows
because $H^2=I$.
\end{proof}

The scope reader satisfies the hypothesis by TS.8, because $\overline H$ is
a projective bijection.

\section{Light and two-level matter}

\subsection{Linear polarization}

Choose a transverse real electric-field amplitude $E\in\RR^2\setminus\{0\}$
at a detector plane. Its unoriented polarization axis is $[E]\in
\RR P^1$. A polarizer reads the scope direction:
\[
E\longmapsto [E].
\tag{TS.20}
\]
The half-turn $E\mapsto-E$ preserves the polarization axis. Wavelength,
detector plane, polarizer convention, intensity threshold, and calibration
fix the measured read.

\subsection{Two-level realization}

Let $\sigma_x,\sigma_y,\sigma_z$ be the Pauli matrices and define
\[
V_\alpha=e^{-i\alpha\sigma_z}.
\tag{TS.21}
\]
For $|+x\rangle$, the Bloch read is
\[
\bigl(\langle\sigma_x\rangle,\langle\sigma_y\rangle\bigr)
=(\cos2\alpha,\sin2\alpha),
\qquad V_{\alpha+\pi}=-V_\alpha.
\tag{TS.22}
\]
The Bloch vector reads the double-cover geometry directly: state phase carries
the local turn and the measured transverse direction carries the scope turn.

\end{document}
